\begingroup
\setlength{\parindent}{0pt}
\setlength{\parskip}{0.45\baselineskip}
\setlength{\abovedisplayskip}{7pt plus 1.5pt minus 1.5pt}
\setlength{\belowdisplayskip}{7pt plus 1.5pt minus 1.5pt}
\setlength{\jot}{5pt}

\section*{How the Proof Works}

The three-dimensional Navier--Stokes regularity problem asks whether a smooth
flow can become singular in finite time. We begin with a smooth,
divergence-free initial velocity \(u_0\in H^\infty(\mathbb R^3)\), of finite
energy, and fix a positive viscosity \(\nu\). The fluid evolves according to
\[
 \partial_tu+(u\!\cdot\nabla)u=-\nabla p+\nu\Delta u,
 \qquad \nabla\cdot u=0,
 \qquad u(0)=u_0.
\]
The hope is simple: viscosity smooths the fluid faster than nonlinear transport can sharpen it.

Strain within the fluid's flow can naturally stretch a vortex. When that
happens, incompressibility narrows the stretched parcel: its volume stays
fixed as its length increases. Stretching can intensify the rotation within
this narrowing region, creating sharper differences between neighboring
velocities. Viscosity smooths those differences. The difficulty is whether
stretching can keep renewing them as the width tends to zero.

We can follow that competition down in scale by rescaling the equation while
keeping viscosity fixed. A contraction of lengths by a factor \(\lambda>1\),
together with an increase in velocity by \(\lambda\), gives
\[
 u_\lambda(x,t)=\lambda u(\lambda x,\lambda^2t),
 \qquad
 p_\lambda(x,t)=\lambda^2p(\lambda x,\lambda^2t).
\]
Every term in the momentum equation gains a factor \(\lambda^3\). Nonlinear
transport and viscous diffusion remain in their original proportion under
this rescaling, however small the scale becomes. The smaller distance over
which viscosity acts is accompanied by faster transport across it. Scaling
alone leaves their competition unresolved.

For kinetic energy, squaring the rescaled velocity contributes
\(\lambda^2\), while the change of spatial variables contributes
\(\lambda^{-3}\). The resulting kinetic energy is smaller by a factor
\(\lambda^{-1}\), even as the velocity increases. To measure the
sharper spatial variation, we use the Sobolev order \(s\). It counts spatial
derivatives in a square-integrable sense, including fractional orders;
increasing \(s\) gives finer variations more weight. Under the rescaling,
\[
 \|u_\lambda(\cdot,t)\|_{\dot H^s}
 =\lambda^{s-1/2}\|u(\cdot,\lambda^2t)\|_{\dot H^s}.
\]
Below \(s=\tfrac12\), the norm falls as the flow concentrates; above it, the
norm rises. At exactly \(s=\tfrac12\), it is unchanged. This is the critical
height, where the measurement preserves its size through the fixed-viscosity
comparison. Invariance under rescaling is a statement about comparing scales;
the critical norm can still change as a solution evolves. Energy lies below
this height, at \(s=0\). Its decrease is compatible with increasingly rapid
motion concentrated into an increasingly small region.

There is also less time between successive stages of concentration. A motion
with length scale \(\ell\) and speed \(U\) has a crossing time \(\ell/U\).
Contracting \(\ell\) by \(\lambda\) while increasing \(U\) by \(\lambda\)
shortens that time by \(\lambda^2\), just as in the rescaled equation. If the
fluid repeated this process at every stage, the durations would satisfy
\[
 \tau_j=\lambda^{-2j}\tau_0,
 \qquad
 \sum_{j=0}^{\infty}\tau_j
 =\frac{\tau_0}{1-\lambda^{-2}}<\infty.
\]
Infinitely many stages could fit into a finite interval, with the width
vanishing and the speed becoming unbounded. This is the hypothetical Zeno
cascade. It explains how a finite-energy fluid could approach a singularity
without running out of time for the required acceleration. Whether the
Navier--Stokes evolution can actually carry out such a cascade still has to
be decided.

In order for a singularity to happen, velocity has to blow up. In order for
that to happen, its acceleration would also have to blow up, which means its
jerk has to blow up, and so on, and so on. Here blow-up means unbounded spatial
supremum norm as a finite terminal time is approached, and acceleration and
jerk refer to the local time derivatives \(\partial_tu\) and
\(\partial_t^2u\). Following a moving fluid particle instead gives the
material acceleration \(\partial_tu+(u\!\cdot\nabla)u\). The necessity in
the local derivative sequence follows directly by integration: for every
finite \(j\), and smooth times \(t_0<t\),
\[
 \|\partial_t^ju(t)\|_\infty
 \le \|\partial_t^ju(t_0)\|_\infty
   +\int_{t_0}^{t}\|\partial_t^{j+1}u(\tau)\|_\infty\,d\tau.
\]
An unbounded left side requires an infinite accumulated size of the next
time derivative, at every step in the escalation.

The equation expresses \(\partial_tu\) through \(\nu\Delta u\), nonlinear transport and
pressure. Differentiating again brings in \(\nu\Delta\partial_tu\); substituting
the equation into that term reaches \(\nu^2\Delta^2u\), along with the
differentiated nonlinear and pressure terms. Repeating the operation keeps
exposing higher spatial derivatives inside the higher time derivatives.

We can see that content at the critical height already found. Write
\(\Lambda=(-\Delta)^{1/2}\), let
\(E_s(t)=\tfrac12\|\Lambda^su(t)\|_2^2\), and consider the scalar quantity
\(H(t)=E_{1/2}(t)\). These are measurements of the evolving velocity field;
their time derivatives combine derivatives of that field. Pairing the
momentum equation with the velocity at Sobolev order \(s\), the viscous term
integrates by parts to \(-\nu\|\Lambda^{s+1}u\|_2^2\). If we write
\[
 \mathcal P_s
 =-\left\langle\Lambda^su,
   \Lambda^s\bigl((u\!\cdot\nabla)u+\nabla p\bigr)\right\rangle_{L^2}
\]
for the transport and pressure contribution, the energy balance is
\(E_s'=\mathcal P_s-2\nu E_{s+1}\). Differentiating the balance at the
critical height, then using it at the newly appearing height, gives
\[
 \begin{aligned}
 H'&=\mathcal P_{1/2}-2\nu E_{3/2},\\
 H''&=\mathcal P_{1/2}'-2\nu\mathcal P_{3/2}
                  +(2\nu)^2E_{5/2}.
 \end{aligned}
\]
At temporal order \(n\), this repeated substitution gives
\((-2\nu)^nE_{n+1/2}\) alongside differentiated transport and pressure terms.
Controlling that complete balance would control the velocity at Sobolev order
\(n+\tfrac12\). Each additional time differentiation reaches a higher spatial
regularity requirement within the evolving fluid. The identity holds while
the flow is smooth; the needed bounds must persist as the endpoint is
approached.

At one fixed Sobolev order, we can control only a finite number of continuous
spatial derivatives. But there is no final order in this construction. Given
any finite derivative order \(k\), we can choose \(n>k+1\), so that
\(n+\tfrac12>k+\tfrac32\), and use Sobolev embedding to control that derivative.
If a common velocity state belongs to every exposed height, with its
finite-energy component retained, it belongs to
\[
 L^2\cap\bigcap_{n=0}^{\infty}\dot H^{n+1/2}
 =\bigcap_{m=0}^{\infty}H^m
 =H^\infty
 \hookrightarrow C^\infty.
\]
For each requested derivative, a sufficiently high finite order does the
work. Smoothness follows because the state satisfies all of those
requirements together. Their bounds can grow with derivative order; smoothness
does not require one number to bound them all.

The temporal side has this interpretation too. At a fixed spatial order
\(m\), controlling every finite time derivative makes the velocity smooth as
an \(H^m\)-valued function of time. Allowing \(m\) to range through every
finite spatial order then accounts for its mixed time and space derivatives.
To use this connection for
continuation, those requirements must hold through the full approach to the
proposed terminal time. Their intersection can then supply a common terminal
state with every spatial derivative, while the time regularity makes the
approach to that state continuous. The fluid would arrive at the supposed
singular time smooth enough to start again.

If a singularity were to form, we would expect to see the opposite behavior:
the velocity and its successive time derivatives escalating while spatial
regularity fails at the limiting time. In the concentrating picture, the
higher Sobolev norms would grow even though the kinetic energy remained
bounded. Every earlier time could still look smooth. The failure would appear
as the loss of control when the endpoint is approached, with bounds valid on
shorter intervals becoming arbitrarily large. Excluding that
loss requires control for every chosen finite set of time and space
derivatives throughout the approach to the endpoint. We have to obtain that
control from the evolution of the given initial velocity.

The initial velocity also determines its initial time derivatives, through
the momentum equation and its differentiated forms:
\begin{flalign*}
 &\partial_tu=\nu\Delta u-(u\!\cdot\nabla)u-\nabla p,&&\\
 &\partial_t^2u=\nu\Delta\partial_tu
   -(\partial_tu\!\cdot\nabla)u
   -(u\!\cdot\nabla)\partial_tu-\nabla\partial_tp,&&\\
 &\partial_t^3u=\nu\Delta\partial_t^2u
   -(\partial_t^2u\!\cdot\nabla)u
   -2(\partial_tu\!\cdot\nabla)\partial_tu
   -(u\!\cdot\nabla)\partial_t^2u-\nabla\partial_t^2p.&&
\end{flalign*}
Each differentiation distributes itself over the two factors in transport.
The coefficients count the ways this can happen. Writing
\(V_r=\partial_t^ru\) and \(Q_r=\partial_t^rp\), the pattern becomes
\[
 V_{r+1}=\nu\Delta V_r
 -\sum_{a=0}^{r}\binom ra(V_a\!\cdot\nabla)V_{r-a}-\nabla Q_r.
\]
Pressure is determined by incompressibility. Taking the divergence of the
original equation gives \(-\Delta p=\partial_i\partial_j(u_i u_j)\).
With the decay normalization, its differentiated form is
\[
 Q_r=R_iR_j\sum_{a=0}^{r}\binom raV_{a,i}V_{r-a,j},
 \qquad R_i=\partial_i(-\Delta)^{-1/2},
\]
where repeated component indices are summed. Starting with \(V_0(0)=u_0\),
these relations determine \(Q_0(0)\), then \(V_1(0)\), then \(Q_1(0)\), and
continue through any finite order. Taking spatial derivatives as well gives
the mixed velocity and pressure derivatives generated by the datum.

A finite piece of this structure retains time derivatives through a chosen
order and measures their spatial regularity through another chosen order.
We call it a rectangle because these two depths, \(N\) and \(M\), can be
chosen independently. The datum-generated rectangle theorem,
Theorem~\ref{thm:datum-rectangles}, establishes the control needed on the
entire maximal interval \([0,T_*)\), under the assumption \(T_*<\infty\):
for every finite \(n,M\),
\[
 V_n\in L^2(0,T_*;H^{M+1}),
 \qquad
 V_{n+1}\in L^2(0,T_*;H^{M-1}).
\]
These adjacent spatial orders reflect the two derivatives in the Laplacian
linking a temporal response to the next. The theorem supplies their
integrability all the way to \(T_*\). Restricting to any earlier cutoff
\(T<T_*\) gives bounds independent of that cutoff; a fixed rectangle
contains only finitely many such norms. The recurrence determines which
coordinates belong in the rectangle, and the datum-generated theorem gives
the full-interval control. Pressure remains attached through its displayed
formula, which also gives the required finite pressure norms.

Enlarging a rectangle adds derivatives and stronger spatial requirements.
Every coordinate already present remains the corresponding derivative of
\(u\) or \(p\). This agreement on overlap is what lets the finite pieces
combine. If we ask for \(r\) time derivatives measured in \(H^m\), one
sufficiently large rectangle contains all of them, and every larger rectangle
agrees with that choice. Ranging over finite \(r,m\) gives the compatible
intersection \(\mathcal I[u_0]\), whose velocity satisfies
\[
 u\in\bigcap_{r,m=0}^{\infty}H^r(0,T_*;H^m(\mathbb R^3)).
\]
Here membership means that each specified finite set of derivatives has a
finite norm. Compatibility ensures that all these statements describe the
original fluid evolving from \(u_0\).

We can now follow that fluid to its endpoint. Fix a spatial order \(m\).
The zeroth temporal row of the intersection gives
\(u\in L^2(0,T_*;H^{m+2})\). Sobolev multiplication and the pressure formula
place \((u\!\cdot\nabla)u+\nabla p\) in \(L^1(0,T_*;H^m)\); the
Laplacian term \(\nu\Delta u\) belongs there too, by Cauchy--Schwarz in time
on the finite interval. The equation
then gives \(u_t\in L^1(0,T_*;H^m)\). Between two late times,
\[
 \|u(t)-u(s)\|_{H^m}
 \le\int_s^t\|u_t(\tau)\|_{H^m}\,d\tau
 \longrightarrow0
 \qquad(s,t\uparrow T_*).
\]
The remaining time integral shrinks to zero, so the velocity converges in
\(H^m\). This holds at every finite \(m\). The limits agree under the
Sobolev embeddings because they are limits of the original velocity, giving
one divergence-free terminal state
\[
 u_*:=\lim_{t\uparrow T_*}u(t)
 \in\bigcap_{m=0}^{\infty}H^m(\mathbb R^3)
 =H^\infty(\mathbb R^3).
\]

The pressure formula and differentiated momentum equation determine every
endpoint time derivative from \(u_*\). Local existence supplies a solution
starting from this smooth state. Its existence time is stable for nearby
late data \(u(t_0)\), and uniqueness identifies the local solutions on their
overlapping intervals. This continues the original solution past \(T_*\),
with the velocity and pressure derivatives agreeing across the join. A finite
maximal lifespan is impossible, so \(T_*=\infty\).

The completed intersection also gives the finite estimates we want to use.
On any finite interval \([0,T]\), choose a time derivative order \(r\), a
spatial derivative order \(k\), and an exponent \(2\le p\le\infty\).
Choosing a spatial height \(m>k+\tfrac32-\tfrac3p\) gives
\[
 \partial_t^ru\in C([0,T];H^m)
 \hookrightarrow L^\infty(0,T;W^{k,p}).
\]
For example, the \(r=0\), \(m=3\) coordinate controls the velocity gradient
in spatial supremum norm, so
\[
 \int_0^T\|\nabla u(t)\|_\infty\,dt
 \le C T\|u\|_{C([0,T];H^3)}<\infty.
\]
The desired estimate determines which finite part of the intersection to
take. Having constructed that intersection and continued the fluid, we can
derive the finite estimates, and use them.

\endgroup
